\chapter{Synchronization}\label{ch:synchronization}

A G17 program has a different ordering mechanism at each of four scopes (\cref{tab:ordering-by-scope}). Visibility,
ordering and forward progress can hold separately: with a threadgroup barrier where a device fence was needed, an
atomic count across threadgroups came out right while the data it guarded was stale (\cref{sec:barriers-fences}).

\begin{xltabular}{\linewidth}{@{}>{\hsize=0.71\hsize}XX>{\hsize=1.12\hsize}XXX@{}}
\caption{Ordering mechanisms by scope.}\label{tab:ordering-by-scope}\\
\toprule
Scope & Mechanism & Visibility & Ordering & Forward progress \\
\midrule
\endfirsthead
\tablecontinued{5}
\toprule
Scope & Mechanism & Visibility & Ordering & Forward progress \\
\midrule
\endhead
one simdgroup & scoreboard slot + consumer wait mask & the register holds the value only for a consumer that waits on its slot & program order on one scoreboard; same-address stores: lowest active lane
wins & none needed (one instruction stream) \\*
simdgroups of one threadgroup & \op{447} at the matching scope & threadgroup memory: any scope; device memory: see 
\cref{sec:barriers-fences} & across the barrier & barriers among a threadgroup's simdgroups completed in
every measured case \\*
threadgroups of one dispatch & \op{14156} on writer and reader, plus an atomic count & measured for the last-threadgroup pattern & from the atomic; a threadgroup barrier gives none & observed only among resident threadgroups; preemption not observed \\*
dispatches & encoder barrier token & measured & serial encoder: after every dispatch & in order (observed) \\*
\bottomrule
\end{xltabular}

\leadhead{Notation.}

\begin{itemize}
\item \emph{Slot}: one of eight scoreboard entries, 0--7.
\item \emph{Fill slot}: the slot a producer names, in operand 1 bits 20--23.
\item \emph{Wait mask}: operand 1 bits 24--31; bit 24 + s waits on slot s.
\item \emph{Late value}: a result that arrives after its producer issues and needs a waiting consumer (\cref{sec:scoreboard}).
\item \emph{Scope immediate}: operand 1 of \op{447} (20, 154, 276, 532, 1144; \cref{sec:barriers-fences}).
\end{itemize}

\section{The scoreboard}\label{sec:scoreboard}

A late-value producer names its fill slot in its own control word (bits 20--23; measured, MM~\mm{0.8}, MM~\mm{6}); the decoder prints slot + 1, so this compiler's first group reads 8 and means slot 7. A consumer waits through
its wait mask. Nothing else orders them; there is no separate wait instruction. What earlier work called a "tag ring" and separately a
"load-use wait bit" are this one field.

If the hardware interlocked on the register, the mask would be only a hint; clearing one wait bit tests this. On an MMA whose operands come from loads filling slot 7, clearing its slot-7 wait bit
(byte0{[}3{]}) made it read the fragments before the loads landed and produce zeros, deterministically; for tensor loads, 50
of 50 dispatches across five processes returned zero (MM~\mm{6}). The mask's layout was established causally (43
dispatches) and then by preregistration: for a load filling slot v - 1 the consumer must set bit 24 + (v - 1), 20 of 20
predicted and confirmed. The imageblock store adds a wrong-slot arm (\cref{tab:imageblock-wait-arms}): moving its wait
from slot 7 to slot 1, which nothing fills, failed like clearing it (MM~\mm{25.96}). A set bit therefore waits on its own slot only;
a late value outstanding on another slot does not hold it.

\Cref{fig:scoreboard}(b) shows a consumer without the bit for the producer's slot reading the register early;
\cref{fig:scoreboard}(a) follows the running example's epilogue through its two late values (\cref{sec:one-program-ir}).

% fig:scoreboard: scoreboard timelines. Panel (a) is the running example's epilogue, with
% the control words read from its bytes (low 32 bits of operand 1). Panel (b) is the
% missing-wait experiment of MM 6. Time is schematic: order, not duration, is measured.
\begin{figure}[htbp]
\centering
\begin{tikzpicture}[x=7.9mm, y=-7.5mm, font=\small,
    issue/.style={circle, fill=ink, inner sep=1.6pt},
    stall/.style={line width=3pt, color=apple!45},
    pend/.style={line width=5pt, color=#1!55},
    ins/.style={anchor=east, align=right, font=\small},
    ctl/.style={anchor=west, font=\scriptsize\ttfamily, text=ink!70}]
  \node[anchor=west, font=\small\bfseries] at (-7.4,-1.3) {(a) The running example's epilogue};
  % instruction rows
  \foreach \y/\n/\o in {0/special-register read/14059, 1/plain 32-bit load/12682,
      2/move immediate/11842, 3/fp32 add/998, 4/plain 32-bit store/17229}
    \node[ins] at (-0.3,\y) {\n\ (\opn{\o})};
  \draw[rule] (0,-0.6) -- (0,6.7);
  % slot lanes
  \node[ins, text=ours] at (-0.3,5.5) {slot 0 pending};
  \node[ins, text=native] at (-0.3,6.3) {slot 7 pending};
  \draw[pend=ours] (0.3,5.5) -- (3.2,5.5);
  \draw[pend=native] (3.4,6.3) -- (8.2,6.3);
  % issue points and stalls
  \node[issue] (i0) at (0.3,0) {};
  \draw[stall] (0.9,1) -- (3.2,1);  \node[issue] (i1) at (3.4,1) {};
  \node[issue] (i2) at (4.0,2) {};
  \draw[stall] (4.6,3) -- (8.2,3);  \node[issue] (i3) at (8.4,3) {};
  \node[issue] (i4) at (9.0,4) {};
  % publish and wait arrows
  \draw[->, ours, line width=0.6pt] (i0) -- (0.3,5.2);
  \draw[->, ours, line width=0.6pt, dashed] (3.2,5.5) -- (3.2,1.2);
  \draw[->, native, line width=0.6pt] (i1) -- (3.4,6.0);
  \draw[->, native, line width=0.6pt, dashed] (8.2,6.3) -- (8.2,3.2);
  % control words
  \node[ctl] at (0.5,-0.32) {0x00100000 publish slot 0};
  \node[ctl] at (3.6,0.68) {0x01800000 wait 0, publish 7};
  \node[ctl] at (4.2,1.68) {0x00000000};
  \node[ctl] at (8.6,2.68) {0x80000000 wait 7};
  \node[ctl] at (9.2,3.68) {no wait: ALU result};
  \node[font=\scriptsize\itshape, text=apple] at (2.05,1.35) {stalled};
  \node[font=\scriptsize\itshape, text=apple] at (6.4,3.35) {stalled};
  \draw[->, ink!50] (0,7.2) -- (10,7.2) node[right, font=\scriptsize\itshape] {time};

  \begin{scope}[yshift=-77mm]
  \node[anchor=west, font=\small\bfseries] at (-7.4,-1.3) {(b) A missing wait (MM~\mm{6})};
  \node[ins] at (-0.3,0) {\op{12674}};
  \node[ins] at (-0.3,1) {tensor MMA, no accumulator (\opn{5107})};
  \node[ins, text=native] at (-0.3,2.2) {slot 7 pending};
  \draw[rule] (0,-0.6) -- (0,2.8);
  \draw[pend=native] (0.3,2.2) -- (6.0,2.2);
  \node[issue] (j0) at (0.3,0) {};
  \node[issue] (j1) at (0.9,1) {};
  \draw[->, native, line width=0.6pt] (j0) -- (0.3,1.9);
  \node[ctl] at (0.5,-0.32) {publish slot 7};
  \node[ctl] at (1.1,0.68) {wait mask without bit 31: issues at once};
  \node[font=\scriptsize, text=native, align=left, anchor=north west, text width=52mm] at (1.1,1.15)
    {reads A and B before the load writes them; the result is zeros in 50 of 50 dispatches};
  \node[draw=native, cross out, line width=1.1pt, minimum size=7pt, inner sep=0pt] at (0.9,1) {};
  \draw[ink, line width=0.8pt] (6.0,-0.25) -- (6.0,0.25);
  \node[ctl, anchor=south] at (6.0,-0.33) {load writes A, B};
  \draw[->, ink!50] (0,3.2) -- (10,3.2) node[right, font=\scriptsize\itshape] {time};
  \end{scope}
\end{tikzpicture}
\caption{Late values and their waits, as the control words specify them. A dot is an instruction's issue; grey is
an instruction held at issue by its wait mask; a coloured bar is a scoreboard slot that a producer has published and
whose value has not yet arrived. A dashed arrow is a slot's arrival releasing its waiter. The control words in (a) are read
from the running example's bytes (the low 32 bits of operand 1: bits 20–23 publish slot + 1, bits 24–31 are the
wait mask). Panel (b) is the measured missing-wait experiment. Positions on the time axis are schematic; only the
order is measured.}
\label{fig:scoreboard}
\end{figure}


The per-form carriers (general loads, the 128-bit vector load, tensor load, the tensor MMA's
scattered mask bits, the 12-byte ALU wait bits, \op{592}) are tabulated in \cref{sec:scoreboard-publish-wait}. One carrier
is not in that table: in 14-byte stores byte9{[}5{]} is slot 7 in this compiler's encoder (\code{agxforge/g17/cc.py}). For
\op{12674}, MM~\mm{6} called the whole field bits 20--27, which includes the load's own wait mask; the fill slot is
bits 20--23 alone. The general-load field is measured through executed consumers
(MM~\mm{25.115}, MM~\mm{25.121}; \code{tools/g17emu.py} \code{load\_hazards}), the staging-file slot only through the decoder (MM~\mm{25.120.1}).

A decoder census gives the rule behind these. For 126 forms across 6,594 Apple programs, every instance's wait bits equal
exactly the slots its source registers are still pending on (\code{isa/g17-wait-mask-census.json}, MM~\mm{25.96}). It covers
stores: the \op{17256} uses every bit 24--31 (13,140 instances in a 314-object census, each one-hot
value present), which refuted an earlier "three-bit store limit" that was an encoder-table artifact (MM~\mm{6}). Tensor stores
are a separate class: all 876 \op{17257} and 118 \op{17258} instances carry a zero wait
byte, and the same-simdgroup MMA-to-tensor-store path is exact with no wait (measured, MM~\mm{6}).

The 8-byte \op{17229} also carries wait bits (15 corpus instances match their pending
slots, \code{isa/g17-wait-mask-census.json}; Apple's \code{tex2d-read} carries 0x1000010 on one, MM~\mm{25.146}). The "no wait bit" note
in \code{agxforge/g17/cc.py} describes this compiler's encoder, which has no carrier for them (\cref{app:a}).

Whether a consumer needs a wait depends on its producer:

\begin{itemize}
\item \emph{Not needed:} ordinary ALU to ALU. 52 of 52 ordered adjacent pairs of ordinary ALU instructions read the previous
  result bit-exactly with no wait (MM~\mm{25.90}, MM~\mm{6}).
\item \emph{Needed:} loads, imageblock reads, vector-load lanes, the packed half2 load, atomic returns, texture fetches and staging
  writes, for every consumer kind. An indirect load whose index is a loaded value read index 0 on every
  lane (1 of 32 right) until its operand 1 also named slot 7 (32 of 32; MM~\mm{25.104}, MM~\mm{25.90}). Shuffles and
  broadcasts of a loaded value read it early (\cref{sec:cross-lane-operations-consumers}). \op{11190}, \op{9986},
  \op{14047} and \op{590} do not interlock on a load (MM~\mm{25.133}), and a store takes its value register at issue
  (MM~\mm{25.96}).
\item \emph{Open:} whether an ALU consumer of an MMA result is interlocked. One ALU consumer placed directly after an MMA carried
  no wait and still read the result; no measured pair had an MMA producer (MM~\mm{6}, MM~\mm{25.90}).
\end{itemize}

Special-register reads fall between these. \op{14059} publishes a slot and Apple's compiler always waits on it, but a
slot-0 read is readable at distance 0 with no wait, and the wait has no measurable cost (measured, MM~\mm{25.117},
MM~\mm{25.121}; probe and counts in \cref{sec:read-latency-whether}).

An unwaited consumer could read the late value, if it arrived in time, or the register's previous contents. To tell these apart the producer's destination was pre-set. With a poison in it, the
four non-interlocking unary forms returned op(poison) on every lane; with a fresh register they returned op(0); a
negative control, an add with its load-wait bit cleared, read the poison on 1,024 of 1,024 lanes, so the instrument can
see a race (measured, MM~\mm{25.133}). The unwaited consumer reads the previous contents. Most observed failures
therefore read 0, and detecting them takes a sentinel or poison value; a zero-filled buffer cannot show them.

An atomic return published on the wrong slot was first found in decoded bytes and then confirmed on hardware
(MM~\mm{25.115}, \code{docs/g17-first-dispatch-trials.md}). This compiler's residency probe counts arrivals with \op{10094} and reads the returned old value. At S = 256 simdgroups every arrival read 1 (a returned 0, plus
one) while the departure maximum was 256. Two explanations fit: simultaneous same-address atomics combine and each requester sees the pre-combine
value (first recorded, later withdrawn), or the consumer read the register before the return arrived. The decoded bytes favoured the second. Marking the atomic's result late
gave its consumer, \op{10279}, this compiler's load-use wait, mask 0b10000000, slot 7, because this
compiler's loads fill slot 7. But the atomic's operand 1 bits 20--23 were 1: it filled slot 0, copied from Apple's
witnesses, where Apple's first consumer waits on bit 24 (120 of 120 corpus instances). The consumer waited on a slot the
atomic never touched. At larger S arrivals had agreed with departures, which had been read as the wait working; in fact
the arrivals were staggered enough that each return landed before its consumer read it. The deciding run changed only the slot. Publishing the atomic on slot 7
(the two programs differ in four bytes) gave an arrival maximum of 256, equal to the departure maximum, 3 of 3, while the
slot-0 arm read 1. Combining would not depend on which slot the return is published on.

Below Metal, the per-lane form behaved the same way. A 12-byte \op{10090}
publishing on slot 0 (operand 1 0x80100000) updated all 64 cells correctly and returned 0 to every lane, and it still
returned 0 with 64 independent two-byte moves inserted before the consumer, so added distance does not help. Changing two atomic bytes to publish on slot 7 (operand 1
0x80800000) returned all 64 old values.

Past eight groups the slot sequence Apple's compiler assigns is scheduler-chosen (twelve loads got 8, 1, 2, 3, 4, 4, 4,
4, 7, 7, 6, 5 in decoded values), so an emitter should use an explicit rule, such as K-step k in slot k mod 8; this
compiler's K = 256 lowerings use slot k mod 7, exact in 3 of 3 cases (MM~\mm{6}, \code{tools/tensorops-model/scoreboard-store-census.json}).
Reusing a slot after its named consumer has waited is safe (measured); reuse before that consumer is unmeasured and
forbidden (MM~\mm{0.14}, MM~\mm{6}). Loads and atomic returns work on slot 7 (measured), but a special-register read
publishing on slot 7 stored nothing, waited or not (measured, MM~\mm{25.121}; \cref{sec:read-latency-whether}).

The glossary's "the return fills scoreboard slot 7" for the uniform-address atomic and \op{11765} describes this
compiler's emission. Apple's corpus fills slots 0 and 1, and every first consumer waits on exactly that slot (25 of 25
for the uniform atomic, MM~\mm{25.115}).

Release acts per lane (\cref{ch:register-file}). A store naming one register as
both value and index, with the index slot releasing it, stores 0 (measured, MM~\mm{25.117}), and an instruction no lane
executes releases nothing (MM~\mm{25.95}).

\section{Barriers and fences}\label{sec:barriers-fences}

\op{447} is one six-byte instruction with scheduling class 3, MayStore, NotDuplicable and side effects.
Only operand 1, the scope immediate, changes with the Metal barrier (\cref{tab:barrier-scopes}; Apple's compiler emits,
\code{isa/g17-barrier-scopes.txt}).

\begin{longtable}{@{}lll@{}}
\caption{Encodings of the Metal threadgroup barriers.}\label{tab:barrier-scopes}\\
\toprule
Metal barrier & Bytes & Scope \\
\midrule
\endfirsthead
\tablecontinued{3}
\toprule
Metal barrier & Bytes & Scope \\
\midrule
\endhead
\code{threadgroup\_barrier(mem\_none)} & \code{075100060000} & 20 \\
\code{threadgroup\_barrier(mem\_device)} & \code{0f6900060000} & 154 \\
\code{threadgroup\_barrier(mem\_threadgroup)} & \code{275100060000} & 276 \\
\code{threadgroup\_barrier(mem\_texture)} & \code{876700060000} & 1144 \\
\code{threadgroup\_barrier(mem\_threadgroup\_imageblock)} & \code{475100060000} & 532 \\
\bottomrule
\end{longtable}

\code{mem\_threadgroup | mem\_device} emits two barriers, one per flag; \code{simdgroup\_barrier} emits nothing
(Apple's compiler emits, \code{isa/g17-atomic-semantics.json} barriers). Scopes 16, 528 and 1138 substituted for 276 each
carried a threadgroup-memory round trip (measured, the threadgroup-barrier notes in \code{tools/g17emu.py}).

Between simdgroups a threadgroup barrier is mandatory (measured, MM~\mm{14}, MM~\mm{17} rule 7). With no barrier, and with
\code{air.simdgroup.barrier} alone (which emits no instruction), a cross-simdgroup hand-off was wrong in 12 of 12 trials and
the consumer read the buffer's pre-write state. Every threadgroup-barrier variant
(execution-only, threadgroup memory, device memory, both) was correct in 12 of 12. A threadgroup store, barrier and load
also exchanged values below Metal, between the two simdgroups of a 64-thread group and in a half-full 32-thread second
group (measured, \code{docs/g17-first-dispatch-trials.md}).

For device memory between simdgroups, MM~\mm{14} observed that the execution-only barrier (scope 20) sufficed, and says so as "an observation about this chip, not a guarantee". The
emulator's race check (MM~\mm{25.197}) orders threadgroup memory across every threadgroup barrier but device memory only across scope 154,
and no measurement requires that restriction. Scope 276 (or any threadgroup barrier scope) is safe for threadgroup
memory between simdgroups, and scope 154 for device memory between simdgroups of one threadgroup. Scope 20 is not
established as safe for memory, and no threadgroup barrier orders anything across threadgroups (\cref{sec:ordering-across-lanes}).

A last-threadgroup pattern compared a threadgroup barrier with a device fence on hardware (MM~\mm{25.140.5}). Every threadgroup
writes a marker, crosses a barrier, and counts itself with \op{10094} on its rank-0
written buffer; the last threadgroup to arrive sums all markers. With the barrier at scope 154 the count was a perfect
permutation, yet the last threadgroup read stale markers: it summed 1,560 of 2,080 at 64 threadgroups and 123,408 of
131,328 at 512. Apple's compiler, given
\code{atomic\_thread\_fence(mem\_device, seq\_cst, thread\_scope\_device)}, emits \op{14156}, operands (0, 186), six
bytes \code{0f2b00060000}, no registers, and places one after the stores and one again inside the last threadgroup's region;
the loads stay ordinary. With those two fences both sums were exact. This compiler emits the same bytes (checked with
Apple's decoder). A fused attention merge built on the pattern was bit-exact at q0 = 0, 5, 128 and 271 over five
rounds each, with NaN-poisoned rows (costs in \cref{ch:performance}).

The imageblock barrier is the threadgroup barrier at scope 532 (\cref{sec:imageblock-per-core-tile}). The emulator admits it only for one simdgroup in
lockstep; more than one simdgroup was never witnessed (MM~\mm{25.145.2}). A barrier inside an exec-mask region is not modelled
there either (\code{tools/g17emu.py}).

Encoder-level barriers come from a static read of the AGX user driver (decoder, MM~\mm{25.141.20}); their costs are measured
in \cref{ch:performance}. A serial compute encoder emits a barrier token after every dispatch. A concurrent encoder emits
one only where \code{memoryBarrierWithScope:} asks; in a serial encoder that call only posts a notification. For a chain
where every dispatch depends on the last, both emit the same tokens. The token is one 32-bit command word, 0x60000160, OR
0x1000 when a pending flag is set; with its boolean arguments it ORs 0x400 or appends 0x60000502, both of unknown
function. The GPU drains at the token.

Hazard tracking never enters the per-dispatch path, since tracked and untracked buffers encode identically inside an
encoder. Consecutive compute encoders in one command buffer are merged, with an unconditional barrier after a concurrent
one. Tracked buffers with one serial encoder per token ran 0.2--0.3~ms per token faster than untracked buffers ordered by
an \code{MTLSharedEvent} between command buffers (\cref{sec:dispatch-encoder-command-buffer}), and untracked buffers
without ordering race (MM~\mm{0.15}).

Two dependent Submits on one queue below Metal, the second reading the first's output, computed exactly, also with a
program switch between them; and one Submit carrying two or three command records, each running a per-lane \code{+7}
read-modify-write on the same words, left 14 and 21 in every lane, so each record saw the previous record's result
(measured, \code{docs/g17-first-dispatch-trials.md}, "Dependent two-stage pure GPU pipeline" and "One to three dependent GPU
computations in one Submit").

\section{Atomics}\label{sec:atomics}

In the decoder (\code{isa/g17-atomic-family.toml}), 208 opcodes share the scheduling class of \op{10094} and flags
(classes 286 and 287 with flags 0x402002400 for device scope, class 329 with 0x402002800 for threadgroup scope) and one
operand skeleton, \code{[dest] imm imm [address] imm [index] imm imm imm [value] imm}, with parts present or absent. Five occur
in the corpus; 29 are reachable from one encoding of it by bit flips. The other 179 were never produced by single-bit
walks to saturation, 132,240 two-bit mutations, 1,365,504 three-bit mutations or 800,000 random encodings. Those
searches are bounded, so the 179 are not shown to be unreachable. Among the family's structural bits, byte4{[}0{]} returns old (to \op{10022}), byte4{[}1{]} selects a threadgroup GPR16 address over a device register pair (\op{11769}, no glossary entry),
byte5{[}2{]} selects an index register (\op{10090}), and byte6{[}2{]} leaves the family.

\begin{xltabular}{\linewidth}{@{}>{\hsize=1.18\hsize}Xl>{\hsize=0.85\hsize}XX@{}}
\caption{Atomic forms by address kind.}\label{tab:atomic-forms}\\
\toprule
Glossary name & Opcode & Address & Standing \\
\midrule
\endfirsthead
\tablecontinued{4}
\toprule
Glossary name & Opcode & Address & Standing \\
\midrule
\endhead
uniform-address atomic, returning old & \opn{10094} & device, uniform, byte offset in operand 6 & measured (below Metal, all operations below; MM~\mm{25.140.3}) \\
uniform-address compare-exchange, returning & \opn{10095} & device, uniform & measured (below Metal) \\
uniform-address no-return atomic, 32-bit & \opn{10022} & device, uniform & Apple's compiler emits \\
uniform-address no-return compare-exchange & \opn{10023} & device, uniform & Apple's compiler emits \\
per-lane-address atomic, returning & \opn{10090} & device, base + index register & measured (below Metal) \\
per-lane-address compare-exchange, returning & \opn{10091} & device, base + index register & measured (below Metal) \\
per-lane-address no-return atomic & \opn{10019} & device, per lane & Apple's compiler emits \\
64-bit no-return atomic, max/min & \opn{9996} & device & Apple's compiler emits \\
threadgroup atomic, returning old & \opn{11765} & threadgroup & measured (MM~\mm{25.115}) \\
threadgroup atomic, no return & \opn{11701} & threadgroup & measured (glossary) \\
\bottomrule
\end{xltabular}

Below, "the uniform atomic" means \op{10094} and "the per-lane atomic"
\op{10090}. "Uniform-address" refers only to the address. The uniform atomic
executes once per active lane, so 32 lanes add 32 and receive 32 distinct old values (measured, MM~\mm{25.140.3}). Its operand 6 is the address's byte offset: values 4, 8 and 16
moved the cell by 1, 2 and 4 words. Earlier claims called operand 6 a source lifetime ("16 writes nothing");
the probe had watched the wrong word (MM~\mm{25.140.3}).

The operation is a four-bit field, named by the decoder's operation operand. The decoder prints operand 2 as
0x40200 + code, so each executed operation's operand is 262656 + code. For the uniform atomic the four bits
are byte4{[}5{]}, byte5{[}3{]}, byte6{[}3{]} and byte7{[}7{]}; \cref{tab:atomic-operation-codes} reproduces all ten encodings Apple's compiler emits for
\code{atomic\_fetch\_\{add,sub,and,or,min,max\}} (both signednesses), \code{atomic\_exchange} and \code{atomic\_fetch\_xor} (decoder plus
Apple's compiler, \code{isa/g17-atomic-family.toml}) and gives the uniform atomic's shortest length for each code.

\begin{xltabular}{\linewidth}{@{}l>{\hsize=1.47\hsize}X>{\hsize=0.74\hsize}X>{\hsize=0.79\hsize}X@{}}
\caption{Atomic operation codes.}\label{tab:atomic-operation-codes}\\
\toprule
Code (operand) & Operation & Shortest length, uniform form & Executed below Metal, uniform form \\
\midrule
\endfirsthead
\tablecontinued{4}
\toprule
Code (operand) & Operation & Shortest length, uniform form & Executed below Metal, uniform form \\
\midrule
\endhead
0 (262656) & add & 10 & yes, returns consumed \\
1 (262657) & sub & 12 & yes \\
2 (262658) & exchange & 12 & yes \\
3 (262659) & compare-exchange: selects \op{10095} & 10 & yes \\
4 (262660) & unsigned min & 12 & yes \\
5 (262661) & signed min & 10 & yes \\
6 (262662) & unsigned max & 12 & yes \\
7 (262663) & signed max & 10 & yes \\
8 (262664) & and & 10 & yes \\
9 (262665) & or & 10 & yes \\
10 (262666) & xor & 12 & yes \\
12 (262668) & fp32 add, emitted by no tested Metal builtin & 10 & yes \\
\bottomrule
\end{xltabular}

Codes 11, 13, 14 and 15 do not decode. The same operand values name the per-lane operations, but the bits carrying them differ
by form and length: in the 12-byte per-lane atomic the fourth bit is byte11{[}4{]}, and a decoder sweep of that form maps the
assembler's three-bit operation with byte11{[}4{]} set and clear onto these codes (\code{docs/g17-first-dispatch-trials.md},
"Per-lane AND old-value return"). An authored "OR" with byte11{[}4{]} set executed exchange, updating every cell to the new
value; clearing that one bit gave operand 262665 and OR (measured). The "eight unassigned raw operation
codes" of \code{isa/g17-atomic-semantics.json} were read off one witness's field placement and are retracted
(\code{isa/g17-atomic-family.toml}; \cref{app:a}).

Each case below ran below Metal, in relaxed order within one dispatch, in two fresh processes with output guards
(\code{docs/g17-first-dispatch-trials.md}; \code{isa/g17-atomic-family.toml}).

\begin{itemize}
\item 64 lanes adding to one cell left exactly the sum (5 to 69 with +1; 5 to 2085 with lane i adding i + 1), and
  every lane received a distinct old value; with lane-varying addends the returned values form a unique chain through all
  64 lanes. Two simdgroups' elected lanes on one uniform atomic received non-overlapping intervals (5--36 and 37--68).
\item Integer operations matched their Metal definitions, including 32-bit modular sub (5 - 32 - 32 = 4294967237), signed against unsigned
  extrema on 0xFFFFFFFF and 0x80000000, and exchange.
\item Compare-exchange takes the register tuple (desired, expected), established from the compiler's witness; matched and
  failed compares both returned the old value.
\item The fp32 add (code 12), uniform and per-lane, returned 1.0f and 2.0f for 1.0f + 1.0f + 1.0f, leaving 3.0f; 64
  per-lane updates of 1.0f from 1.0f reached 65.0f and returned 1.0f through 64.0f once each. Each update rounds to
  nearest even: half-ULP ties left an even mantissa unchanged and moved an odd one up once (0x3F800001 to 0x3F800002) and
  no further. Subnormal inputs, operands and results flush to zero. A quiet NaN seed 0x7FC12345 was returned raw to its
  first caller and then stored as canonical 0x7FC00000; $-0$ plus $+0$ returned $-0$ and stored $+0$; infinities were preserved.
\end{itemize}

Atomics have been run only on the rank-0 buffer, through Metal and below it; this compiler refuses an
atomic whose identity base and binding rank disagree (\code{agxforge/g17/cc.py}; MM~\mm{25.140.3}). Threadgroup atomics need a
threadgroup region, which only an ordinary threadgroup load or store provisions (\cref{sec:threadgroup-memory}).

Atomic return values are late; the fill slot is operand 1 bits 20--23 minus 1, and slot choice is in \cref{sec:scoreboard}. The 10-byte uniform atomic also releases its value: operand 9 is the value's lifetime and only 16 is
representable, so a second atomic sharing a constant read 0 (A{[}0{]} = 0 against 24; measured by this compiler's fuzzer,
MM~\mm{25.144.6}, \code{agxforge/g17/cc.py}); a value read again must be copied first.

Validation with the return value consumed:

\begin{itemize}
\item The uniform atomic was validated below Metal for all eleven of its operations in \cref{tab:atomic-operation-codes} (every code except 3), each under two-group
  contention with the returned old value consumed; \op{10095} at match and
  mismatch. Through Metal it was validated for add, in the residency probe (MM~\mm{25.115})
  and the compiler's wave-broadcast program (\code{agxforge/g17/cc.py}).
\item The per-lane atomic ran below Metal update-only for add, and, or, sub, xor, signed and unsigned min and max, exchange and
  fp32 add; returns through slot 7 for add (also contended, both fixed and lane-varying addends), and, or, sub, xor and
  fp32 add (also contended, half-ULP and NaN cases). Compiler-generated 12-byte programs for add, and, or, sub and xor
  passed with a direct indexed-store consumer, plus an ordinary add consumer for add and xor.
\item \Op{10091} passed with a hand-authored slot-7 return (one matching and 63 failing compares), then with
  compiler-generated programs at index placements R0, R1, R5, R7 and R12 (ten Submits); other placements are decode-checked
  predictions, and placement 13 is refused.
\item \Op{11765} was dispatched through Metal with its return consumed; its \op{11667} consumer had read both returns with an empty wait mask until fixed, where Apple waits 10 of 10 (MM~\mm{25.115}).
\item Not validated are per-lane min and max returns, unsigned long-form returns, other consumer shapes (refused by this
  compiler), 64-bit atomics (\op{9996}, compiled only), and any atomic on a buffer other
  than rank 0.
\end{itemize}

An atomic does not order other memory (\cref{sec:barriers-fences}); ordering ordinary stores through an atomic needs
\op{14156}. The emulator does not
track atomics in its race check and does not model acquire or release through atomics or the fence, so it refuses any
program that synchronizes threadgroups through an atomic counter (MM~\mm{25.197}).

\section{Ordering by scope}\label{sec:ordering-across-lanes}

\subsection*{Lanes of one simdgroup}

\begin{itemize}
\item Program order holds on one scoreboard: a store followed by a load in one simdgroup needs no fence, as two separately
  generated bodies joined at a shared buffer showed, including a cross-lane stress where each lane read another lane's
  write (measured, MM~\mm{13}). Late values still need the consumer's wait.
\item Same-address stores with \op{17229}: the lowest active lane wins, by lane and not by value (all
  lanes: lane 0; lanes 0--15 under an exec region: lane 0; odd lanes: lane 1; reversed values: lane 0; five runs each;
  measured, MM~\mm{25.145.4}, \code{isa/g17-execution-storewinner-results.json}).
\item Registers never cross a simdgroup: a cross-simdgroup \code{simd\_shuffle} returns the caller's own value, 256 of 256
  (measured, MM~\mm{14}).
\end{itemize}

\subsection*{Simdgroups of one threadgroup}

A threadgroup barrier is the only mechanism at this scope (\cref{sec:barriers-fences}); same-value writes by
several simdgroups are harmless by construction (no schedule changes the bytes), and shipped code relies on that: the fused attention appends a
new K/V row from every simdgroup of every threadgroup and each reads its own copy back (modelling rule, MM~\mm{25.197}).

\subsection*{Threadgroups of one dispatch}

MM~\mm{0.14}'s map says "cross-threadgroup handoff requires a
dispatch boundary". It was written about tensor fragment hand-offs before any in-dispatch mechanism had been measured.
MM~\mm{25.140.5} and MM~\mm{25.141.1} have since measured two, so the dispatch boundary is the safe default, and fenced
atomic patterns also work within one dispatch.

Device stores from one threadgroup are visible to another in the same dispatch after \op{14156} on the writer side and
again on the reader side (measured for the last-threadgroup pattern, MM~\mm{25.140.5}); a threadgroup barrier at any
scope does not make them visible. An atomic count decides which threadgroup is last (atomicity is measured in
\cref{sec:atomics}).

The last-threadgroup pattern needs no forward progress, because no threadgroup waits; whoever counts last does the
follow-on work. Any pattern that waits needs the producer resident. In the forward-progress probe the first resident
wave spun to its cap before the producer, dispatched last, ran at all (measured, MM~\mm{25.141.1}; the grid sizes and
counts are in \cref{sec:forward-progress-in-order}). A grid barrier among resident
threadgroups, fenced and capped, kept totals exact in every run and never reached its cap (measured, MM~\mm{25.141.1}).
MM~\mm{17} rule 40 allows a wait on another threadgroup only when every participant fits in measured resident capacity
at the kernel's own register and threadgroup-memory use, with capped spins and work pulled from an atomic queue. The barrier took
133.1~µs per 32~MB phase against 120.8~µs with dispatch boundaries (\cref{sec:dispatch-encoder-command-buffer}).

The residency of 1,024-thread threadgroups is unresolved. The forward-progress probe implies at least
10 co-resident per core, while the driver formula gives 3 (MM~\mm{25.141.1}, MM~\mm{25.141.20}); no residency count may be assumed, and
a spin-waiting grid of them must not be sized from either number (\cref{sec:residency-occupancy-rules}, rule 5).

\subsection*{Dispatches}

In the tested configurations threadgroups were dispatched in order and no spinning threadgroup gave way to a producer
dispatched after it (MM~\mm{25.141.1}; \cref{sec:forward-progress-in-order}). No guarantee of that behaviour has been
established, so a program cannot depend on preemption. A serial encoder places
a barrier token after every dispatch; a concurrent encoder lets a dependent kernel start early, but its threadgroups queue
behind all of the producer's, so only the producer's last-wave tail overlaps (measured, MM~\mm{25.141.1}; costs \cref{ch:performance}).
A concurrent encoder needs \code{memoryBarrierWithScope:} before every dispatch that reads
an earlier dispatch's output.

A line of a write-combined buffer that the GPU read recently stays stale in the GPU's L2 after the CPU updates it
(measured, MM~\mm{0.15}).

\textbf{Not known.} The cross-simdgroup winner of same-address stores; whether an execution-only barrier orders device
memory by design; threadgroup barriers in partial groups of other sizes and with inactive lanes accessing memory; why a
special-register read on slot 7 stores nothing; signaling NaNs and other NaN payloads in the fp32 atomic add; acquire and
release semantics in general (only the fence-count-fence pattern is measured); the function of the barrier-token
variants 0x400 and 0x60000502; the ordering of other command shapes below Metal, and whether command records can
overlap.

\textbf{Evidence.} MM~\mm{0.8}, MM~\mm{0.14}, MM~\mm{0.15}, MM~\mm{6}, MM~\mm{13}, MM~\mm{14}, MM~\mm{17} rules 4, 7, 35 and 40, MM~\mm{25.90}, MM~\mm{25.95}, MM~\mm{25.96},
MM~\mm{25.104}, MM~\mm{25.115}, MM~\mm{25.117}, MM~\mm{25.120.1}, MM~\mm{25.121}, MM~\mm{25.133}, MM~\mm{25.140.3}, MM~\mm{25.140.5}, MM~\mm{25.141.1}, MM~\mm{25.141.20},
MM~\mm{25.144.6}, MM~\mm{25.145.2}, MM~\mm{25.145.4}, MM~\mm{25.146}, MM~\mm{25.197}; \code{isa/g17-barrier-scopes.txt}, \code{isa/g17-atomic-family.toml},
\code{isa/g17-atomic-semantics.json}, \code{isa/g17-wait-mask-census.json}, \code{isa/g17-execution-storewinner-results.json},
\code{docs/g17-first-dispatch-trials.md}, \code{agxforge/g17/cc.py}, \code{tools/g17emu.py},
\code{tools/tensorops-model/scoreboard-store-census.json}.
